\documentclass[11pt,a4paper]{article}

% =========================================================
% Packages
% =========================================================
\usepackage[a4paper,margin=2.5cm]{geometry}
\usepackage{amsmath,amssymb}
\usepackage{booktabs}
\usepackage{longtable}
\usepackage{array}
\usepackage{multirow}
\usepackage{float}
\usepackage{graphicx}
\usepackage{xcolor}
\usepackage{hyperref}
\usepackage{enumitem}
\usepackage{microtype}
\usepackage{siunitx}

\hypersetup{
    colorlinks=true,
    linkcolor=blue,
    citecolor=blue,
    urlcolor=blue
}

\sisetup{
    output-decimal-marker={.},
    group-separator={,},
    group-minimum-digits=4
}

% =========================================================
% Document information
% =========================================================
\title{\textbf{Statistical and SQL-Based Customer Analytics for\\
E-Commerce Decision Making}}
\author{Mahdiyeh Mirzaei}
\date{\today}

\begin{document}

\maketitle

\begin{abstract}
This study presents a statistical and SQL-based analysis of customer
behavior in an e-commerce dataset containing 60 customers and 26 original
variables. The analysis integrates descriptive statistics, data-quality
diagnostics, non-parametric inference, correlation analysis, customer-value
scoring, churn-risk classification, segmentation, and business-oriented
aggregation. The primary objective is to translate transactional and
customer-level observations into statistically interpretable evidence for
customer relationship management and marketing decisions.

The results indicate substantial heterogeneity in customer value and
purchasing behavior. Total spending was strongly non-normal
(Shapiro--Wilk $W=0.708$, $p<0.001$), while purchase count
($W=0.957$, $p=0.033$), average order value
($W=0.915$, $p<0.001$), age ($W=0.938$, $p=0.004$), and satisfaction
score ($W=0.879$, $p<0.001$) also departed from normality. Spearman
rank correlations showed a moderate positive association between purchase
count and total spending ($\rho=0.613$), and a weaker positive association
between average order value and total spending ($\rho=0.434$).
Customer-value scoring divided the sample into three equally sized groups:
20 high-value, 20 medium-value, and 20 low-value customers. The high-value
segment generated approximately 70.8\% of total revenue.

A churn-risk rule based on more than 180 days since the last purchase
classified 36 of 60 customers (60\%) as at risk. Importantly, 12 customers
were simultaneously classified as high-value and churn-risk, identifying
a strategically critical retention segment. K-means clustering selected
$k=3$ as the preferred solution, with a silhouette score of 0.316; a
Kruskal--Wallis test confirmed statistically significant differences
between clusters ($H=16.109$, $p=0.000318$). The discount analysis did not
provide evidence of a meaningful spending difference between discount and
non-discount customers (Mann--Whitney $p=0.4512$; Cohen's $d\approx-0.022$;
bootstrap 95\% confidence interval for the mean difference:
$[-2146.19,\,2400.88]$).

Geographically, Mashhad produced the highest total revenue
(64,118.44) and the highest mean customer spending (5,828.95), although
its mean satisfaction score was only 2.91. Overall, the findings support
targeted retention, value-based segmentation, geographically informed
marketing, and selective rather than universal discounting.
\end{abstract}

\noindent\textbf{Keywords:} customer analytics; SQL; descriptive statistics;
non-parametric statistics; customer lifetime value; churn risk; K-means;
Mann--Whitney U; Kruskal--Wallis; e-commerce analytics.

% =========================================================
\section{Introduction}

The rapid growth of e-commerce has increased the availability of
customer-level transactional and behavioral data. However, the presence
of data alone does not guarantee effective decision making. Managers need
statistically defensible summaries that distinguish systematic patterns
from random variation and identify customer groups with economically
meaningful differences.

The present study uses SQL outputs as the analytical foundation and
combines database aggregation with statistical reasoning. The analysis
addresses five central business questions:

\begin{enumerate}[label=\arabic*.]
    \item Which customers generate the greatest economic value?
    \item Which customers are potentially vulnerable to churn?
    \item Which geographical markets have the strongest revenue potential?
    \item Is discount usage associated with higher customer spending?
    \item Can customers be separated into statistically and behaviorally
    meaningful segments?
\end{enumerate}

The analysis is deliberately designed around the distributional properties
of the observed variables. Because several important measures violate
normality, rank-based and distribution-free procedures are emphasized
where appropriate.

% =========================================================
\section{Data and Analytical Framework}

The analytical dataset contains 60 customers and 26 original variables.
The principal variables include customer identifier, age, city, province,
membership tier, purchase count, average order value (AOV), total spending,
days since last purchase, discount usage, returned items, and satisfaction
score.

Additional analytical variables were constructed from the raw attributes:

\begin{itemize}
    \item customer value score;
    \item customer value segment;
    \item churn-risk indicator;
    \item churn-risk label;
    \item cluster membership.
\end{itemize}

The SQL layer was used for data validation, aggregation, ranking,
cross-tabulation, and business summaries. Statistical analysis was then
used to assess distributional characteristics, associations, group
differences, and segmentation quality.

% =========================================================
\section{Data Quality Assessment}

The first stage consisted of structural and logical quality checks.
The dataset contained 60 customers and no duplicate records. The principal
analytical variables did not contain material missingness that prevented
the planned analyses.

A consistency check compared observed total spending with the expected
amount implied by purchase count and average order value:

\begin{equation}
E_i = \text{PurchaseCount}_i \times \text{AOV}_i.
\end{equation}

Two observations required additional review. Customer 1030 had observed
total spending of 25,000.00 compared with an expected value of approximately
4,079.14, producing a discrepancy of about 20,920.86. Customer 1040 had
observed spending of 126.00 compared with an expected value of approximately
1,280.44, producing a discrepancy of about 1,154.44.

These observations should be validated against the underlying transaction
system before the data are used for financial reconciliation or automated
decision rules.

Six records were also flagged for review in the returned-items analysis.
Customers requiring review had substantially lower mean spending
(596.10) and lower mean satisfaction (2.50) than customers without the
flag (mean spending 4,040.29; mean satisfaction 3.04). This pattern
suggests that return-related anomalies may also be associated with weaker
customer experience.

% =========================================================
\section{Descriptive Statistical Analysis}

The total observed revenue was 221,752.15, with a mean customer spending
of approximately 3,695.87 and a mean AOV of approximately 213.16.
The mean purchase count was approximately 17.38 transactions per customer.

The distribution of customer spending was highly heterogeneous. The
largest spending values were substantially above the central tendency,
which is consistent with the observed non-normal distribution and
motivates the use of robust and non-parametric procedures.

\subsection{Normality Assessment}

Normality was evaluated using the Shapiro--Wilk test. The results are
summarized in Table~\ref{tab:normality}.

\begin{table}[H]
\centering
\caption{Shapiro--Wilk normality tests}
\label{tab:normality}
\begin{tabular}{lrrl}
\toprule
Variable & $W$ & $p$-value & Interpretation \\
\midrule
Age & 0.938 & 0.004 & Non-normal \\
Purchase Count & 0.957 & 0.033 & Non-normal \\
Average Order Value & 0.915 & $<0.001$ & Non-normal \\
Total Spending & 0.708 & $<0.001$ & Strongly non-normal \\
Last Purchase Days & 0.970 & 0.142 & No evidence against normality \\
Satisfaction Score & 0.879 & $<0.001$ & Non-normal \\
\bottomrule
\end{tabular}
\end{table}

At the conventional $\alpha=0.05$ significance level, normality was
rejected for age, purchase count, AOV, total spending, and satisfaction.
Consequently, rank-based statistics and robust summaries are more
appropriate than relying exclusively on classical parametric assumptions.

% =========================================================
\section{Correlation Analysis}

Because several variables were non-normal and the sample size was limited,
Spearman's rank correlation coefficient was used to quantify monotonic
associations.

\begin{table}[H]
\centering
\caption{Selected Spearman rank correlations}
\label{tab:corr}
\begin{tabular}{lrr}
\toprule
Variable Pair & Spearman $\rho$ & Interpretation \\
\midrule
Purchase Count -- Total Spending & 0.613 & Moderate positive \\
AOV -- Total Spending & 0.434 & Weak-to-moderate positive \\
Last Purchase Days -- Satisfaction & -0.293 & Weak negative \\
Age -- Purchase Count & 0.265 & Weak positive \\
\bottomrule
\end{tabular}
\end{table}

The strongest observed association was between purchase count and total
spending ($\rho=0.613$), indicating that customers who purchase more
frequently tend to generate greater cumulative revenue. Nevertheless,
the relationship is not deterministic, because spending also depends on
order value.

The correlation between AOV and total spending was positive but weaker
($\rho=0.434$), suggesting that both purchase frequency and order size
contribute to customer value.

% =========================================================
\section{Customer Value Scoring}

A composite customer-value score was constructed from standardized
measures of total spending, purchase count, and AOV:

\begin{equation}
CVS =
0.50Z(\text{Total Spending})
+0.30Z(\text{Purchase Count})
+0.20Z(\text{AOV}).
\end{equation}

The weighting scheme assigns the greatest contribution to cumulative
economic value while retaining behavioral information from frequency and
order size.

Customers were divided into three equal-sized quantile groups:
High Value, Medium Value, and Low Value.

\begin{table}[H]
\centering
\caption{Customer-value segments}
\label{tab:value}
\begin{tabular}{lrrrr}
\toprule
Segment & $n$ & Mean Spending & Total Revenue & Revenue Share \\
\midrule
High Value & 20 & 7,846.92 & 156,938.43 & 70.77\% \\
Medium Value & 20 & 2,278.88 & 45,577.59 & 20.55\% \\
Low Value & 20 & 961.81 & 19,236.13 & 8.68\% \\
\bottomrule
\end{tabular}
\end{table}

The concentration of approximately 70.77\% of revenue within one-third
of the customer base is a major managerial finding. It indicates a
high degree of revenue concentration and supports differentiated
customer relationship strategies.

% =========================================================
\section{Churn-Risk Analysis}

Churn risk was operationalized using a recency threshold: customers with
more than 180 days since their last purchase were classified as
churn-risk customers.

\begin{table}[H]
\centering
\caption{Churn-risk comparison}
\label{tab:churn}
\begin{tabular}{lrrrr}
\toprule
Group & $n$ & Mean Spending & Mean Purchases & Mean Last-Purchase Days \\
\midrule
Churn Risk & 36 & 3,816.63 & 16.89 & 265.14 \\
Active & 24 & 3,514.73 & 18.13 & 102.08 \\
\bottomrule
\end{tabular}
\end{table}

The churn-risk group contains 36 of 60 customers, corresponding to a
60\% observed risk rate under the selected rule. The most important
strategic subgroup consists of the 12 customers who are simultaneously
high-value and churn-risk.

This overlap is more informative than the overall churn rate because
customer loss is economically asymmetric: losing a high-value customer
is potentially more consequential than losing a low-value customer.

It should be emphasized that the 180-day rule is a classification
criterion rather than a probabilistic survival model. It therefore
identifies operational risk rather than estimating a causal probability
of future churn.

% =========================================================
\section{Geographical Analysis}

City-level aggregation revealed substantial heterogeneity in revenue.

\begin{table}[H]
\centering
\caption{Selected city-level performance indicators}
\label{tab:city}
\begin{tabular}{lrrrr}
\toprule
City & Customers & Total Revenue & Mean Spending & Mean Satisfaction \\
\midrule
Mashhad & 11 & 64,118.44 & 5,828.95 & 2.91 \\
Tabriz & 12 & 36,192.49 & 3,016.04 & 3.00 \\
Karaj & 7 & 26,189.81 & 3,741.40 & 2.86 \\
Ahvaz & 8 & 23,384.78 & 2,923.10 & 3.50 \\
Isfahan & 5 & 21,975.32 & 4,395.06 & 4.40 \\
Shiraz & 6 & 20,715.92 & 3,452.65 & 2.00 \\
Rasht & 5 & 15,168.54 & 3,033.71 & 2.60 \\
Tehran & 6 & 14,006.85 & 2,334.48 & 2.67 \\
\bottomrule
\end{tabular}
\end{table}

Mashhad generated the highest total revenue and the highest mean
customer spending. However, its satisfaction score of 2.91 is not
particularly strong. Thus, a pure revenue-maximization strategy may
overlook customer-experience risk. Marketing investment in Mashhad
should therefore be accompanied by service-quality monitoring.

Isfahan provides a contrasting profile: its mean satisfaction score of
4.40 is the highest in the dataset, while its customer base is relatively
small. This combination may justify controlled customer-acquisition
experiments.

% =========================================================
\section{Discount Effect Analysis}

The relationship between discount usage and customer spending was
examined using the Mann--Whitney U test because spending is strongly
non-normal.

The mean spending was 3,642.42 for discount users and 3,736.74 for
non-users. The observed difference was not statistically significant:

\begin{equation}
p=0.4512.
\end{equation}

The standardized effect size was approximately

\begin{equation}
d\approx -0.022,
\end{equation}

which is practically negligible. A bootstrap 95\% confidence interval
for the mean spending difference was approximately

\begin{equation}
[-2146.19,\;2400.88].
\end{equation}

Because the confidence interval includes zero and the effect size is
close to zero, the data do not provide evidence that discount usage is
associated with a systematic increase in spending.

This result should not be interpreted as proof that discounts never work.
Rather, within this observational sample, there is insufficient evidence
of a general spending uplift attributable to discount usage. Targeted
experimentation or randomized A/B testing would be preferable for causal
evaluation.

% =========================================================
\section{Membership and Channel Analysis}

Membership tiers exhibited heterogeneous economic profiles. Gold
customers generated the highest total revenue, 75,153.20, while Silver
customers had the highest mean spending per customer, approximately
4,179.93. VIP customers had a higher mean satisfaction score of 3.40,
whereas Silver customers had the lowest satisfaction score at 2.50.

Payment-channel analysis showed that cash payments generated the largest
total revenue (80,781.87), followed by online wallet payments (71,402.63)
and card payments (69,567.65).

Device analysis showed that web users generated approximately 78,876.06
in revenue, followed by Android users at 74,945.08 and iPhone users at
67,931.01. These figures describe observed channel performance and should
not be interpreted causally without controlling for customer composition
and exposure.

% =========================================================
\section{Age and Satisfaction Profiles}

The 55+ age group exhibited the highest mean spending, approximately
6,592.10, and an average purchase count of 21.35. Younger age groups
generally exhibited lower spending levels.

The relationship between satisfaction and spending was not monotonic
enough to support a simple linear interpretation. Customers with a
satisfaction score of 4 had mean spending of approximately 5,448.63,
whereas the score-5 group had mean spending of approximately 2,677.52.

This pattern demonstrates why customer satisfaction should not be treated
as a direct proxy for economic value. Satisfaction and spending represent
different dimensions of customer behavior and should be modeled jointly.

% =========================================================
\section{Cluster Analysis}

K-means clustering was used to identify behavioral customer segments.
The preferred number of clusters was $k=3$, with a silhouette coefficient
of 0.316. Although the silhouette value indicates moderate rather than
strong separation, it provides useful evidence of meaningful
heterogeneity in the customer base.

\begin{table}[H]
\centering
\caption{Cluster profiles}
\label{tab:clusters}
\begin{tabular}{lrrrrrr}
\toprule
Cluster & $n$ & Mean Age & Purchases & AOV & Spending & Satisfaction \\
\midrule
0 & 6 & 58.50 & 30.33 & 348.60 & 14,188.76 & 3.17 \\
1 & 23 & 43.96 & 8.30 & 309.07 & 2,603.76 & 2.65 \\
2 & 31 & 40.65 & 21.61 & 115.78 & 2,475.26 & 3.19 \\
\bottomrule
\end{tabular}
\end{table}

Cluster 0 is particularly important because only six customers generate
an average spending of 14,188.76. The total revenue attributed to this
cluster is approximately 85,132.55, representing 38.39\% of cluster-level
revenue.

To test whether the clusters differ statistically, the Kruskal--Wallis
test was applied. The result was

\begin{equation}
H=16.109,\qquad p=0.000318.
\end{equation}

At $\alpha=0.05$, the null hypothesis of identical distributions across
clusters is rejected. Thus, the segmentation is supported by statistically
significant between-cluster differences, although the magnitude and
business meaning of each cluster should still be interpreted substantively.

% =========================================================
\section{Machine-Learning-Based Churn Diagnostics}

As a supplementary predictive analysis, Logistic Regression and Random
Forest models were evaluated using stratified five-fold cross-validation.
The predictor set included age, purchase count, AOV, total spending,
satisfaction score, and returned items. Last-purchase days was excluded
from the predictive feature set to reduce target leakage because the
churn label itself is constructed from recency.

The mean ROC-AUC for Logistic Regression was 0.564
($SD=0.084$), while Random Forest achieved 0.478
($SD=0.096$). These values indicate weak discrimination and suggest that
the available behavioral variables do not provide a strong predictive
signal for the operational churn label.

Random Forest feature importance ranked total spending (0.235), AOV
(0.214), satisfaction (0.167), age (0.156), purchase count (0.149), and
returned items (0.078) as the principal predictors.

The relatively weak cross-validated AUC values imply that the current
dataset is better suited to descriptive segmentation and risk
prioritization than to high-confidence individual churn prediction.

% =========================================================
\section{Integrated Statistical Interpretation}

The statistical evidence supports four broad conclusions.

First, customer value is highly concentrated. The high-value segment
contains only one-third of customers but generates approximately 70.77\%
of revenue. This concentration increases the economic importance of
retention among high-value customers.

Second, the customer base exhibits substantial behavioral heterogeneity.
The non-normal distributions, moderate correlations, and significant
cluster differences all indicate that a single average customer profile
would be an inadequate representation of the population.

Third, churn risk is operationally substantial. Sixty percent of the
sample meets the selected recency-based risk criterion, and 12 customers
combine high economic value with elevated recency risk. These customers
should be prioritized before broad, undifferentiated retention campaigns.

Fourth, the evidence does not support universal discounting as a reliable
revenue-growth mechanism. The Mann--Whitney test, near-zero effect size,
and bootstrap confidence interval are all consistent with the conclusion
that discount usage did not produce a detectable general spending
difference in this sample.

% =========================================================
\section{Managerial Implications}

The statistical findings translate into several actionable strategies:

\begin{enumerate}
    \item \textbf{Prioritize high-value churn-risk customers.}
    The 12 customers in the intersection of the high-value and churn-risk
    groups represent the most economically sensitive retention segment.

    \item \textbf{Use value-based rather than uniform marketing.}
    Since revenue is strongly concentrated, customer-level resource
    allocation should reflect expected economic value.

    \item \textbf{Treat Mashhad as a high-potential but experience-sensitive
    market.} Its revenue performance is strong, but satisfaction should
    be monitored before substantially increasing acquisition expenditure.

    \item \textbf{Use discounts selectively.}
    The lack of a significant spending difference suggests that blanket
    discounting is not justified by the current evidence.

    \item \textbf{Investigate data anomalies before automated decisions.}
    The discrepancies identified for customers 1030 and 1040 may materially
    affect customer-value calculations and should be reconciled.

    \item \textbf{Use cluster-specific campaigns.}
    The statistically significant differences between clusters support
    differentiated communication, loyalty benefits, and retention
    strategies.

    \item \textbf{Improve predictive data collection.}
    The weak cross-validated ROC-AUC values indicate that additional
    temporal, behavioral, and interaction-level features would be useful
    for future churn modeling.
\end{enumerate}

% =========================================================
\section{Limitations}

Several limitations should be considered when interpreting the results.
First, the sample size is small ($n=60$), which limits statistical power
and the stability of predictive models. Second, the analysis is
observational; therefore, associations between discount usage, spending,
satisfaction, and churn risk should not be interpreted as causal effects.

Third, the churn definition is rule-based rather than derived from a
time-to-event survival model. Fourth, the K-means solution provides
moderate rather than strong separation, as reflected by a silhouette
score of 0.316. Finally, data-quality discrepancies in selected records
may influence individual-level rankings and should be resolved before
operational deployment.

% =========================================================
\section{Overall Conclusion}

This study demonstrates how SQL-based customer analytics can be extended
into a statistically grounded decision-support framework. The analysis
of 60 customers revealed pronounced heterogeneity in purchasing behavior,
customer value, geographic performance, and churn risk.

The strongest finding is the concentration of economic value: the
high-value third of the customer base accounts for approximately 70.77\%
of total revenue. At the same time, 36 customers, representing 60\% of
the sample, satisfy the operational churn-risk criterion. The intersection
of these two dimensions identifies 12 high-value customers who are also
at risk of churn. From a managerial perspective, this group represents
the most immediate retention opportunity because the potential economic
cost of losing these customers is substantially greater than that of
losing low-value customers.

The inferential analysis further shows that customer behavior cannot be
adequately summarized by simple means or normal-theory assumptions.
Several variables are significantly non-normal, motivating the use of
Spearman correlation, Mann--Whitney U, and Kruskal--Wallis procedures.
The clustering analysis identifies three statistically distinguishable
customer profiles, while the weak predictive AUC values indicate that
the current feature set is insufficient for reliable individual-level
churn prediction.

The discount analysis is particularly important for managerial
decision-making. No statistically significant spending difference was
detected between discount users and non-users ($p=0.4512$), the estimated
effect size was negligible, and the bootstrap confidence interval
contained zero. Consequently, discounts should be viewed as a tactical
instrument requiring targeting and experimentation rather than as a
general-purpose mechanism for increasing customer spending.

Finally, Mashhad emerged as the strongest revenue market, but its
moderate satisfaction level highlights the need to balance acquisition
with customer-experience improvement. Overall, the evidence supports a
data-driven strategy centered on value-based segmentation, targeted
retention, selective incentives, geographic prioritization, and stronger
data-quality controls. Future research should extend the analysis with
transaction-level temporal data, survival analysis, causal experiments
for discount effectiveness, and richer behavioral features to improve
predictive customer analytics.

% =========================================================
\section*{Key Statistical Indicators}

\begin{table}[H]
\centering
\caption{Summary of principal analytical indicators}
\begin{tabular}{lr}
\toprule
Indicator & Value \\
\midrule
Sample size & 60 customers \\
Total revenue & 221,752.15 \\
Mean customer spending & 3,695.87 \\
Mean AOV & 213.16 \\
High-value customers & 20 \\
High-value revenue share & 70.77\% \\
Churn-risk customers & 36 (60\%) \\
High-value and churn-risk customers & 12 \\
Best city by total revenue & Mashhad \\
Mashhad total revenue & 64,118.44 \\
Mashhad mean spending & 5,828.95 \\
Discount test: Mann--Whitney $p$ & 0.4512 \\
Discount effect size: Cohen's $d$ & $\approx -0.022$ \\
Discount mean-difference bootstrap 95\% CI & $[-2146.19,2400.88]$ \\
Optimal K-means $k$ & 3 \\
Silhouette score & 0.316 \\
Kruskal--Wallis $H$ & 16.109 \\
Kruskal--Wallis $p$ & 0.000318 \\
Logistic Regression mean ROC-AUC & 0.564 \\
Random Forest mean ROC-AUC & 0.478 \\
\bottomrule
\end{tabular}
\end{table}

\end{document}
